\documentclass[11pt]{article}
\usepackage[margin=1in]{geometry}
\usepackage{amsmath,amssymb,amsthm,mathtools}
\usepackage[hidelinks]{hyperref}
\newtheorem{theorem}{Theorem}
\newtheorem{lemma}[theorem]{Lemma}
\newtheorem{corollary}[theorem]{Corollary}
\newcommand{\Q}{\mathbb Q}
\newcommand{\F}{\mathbb F}
\newcommand{\vp}{v_p}
\title{Prime-gap bounds and small residual kernels for factorial square divisors}
\author{User-supplied deductions checked and integrated; no priority claim}
\date{16 November 2026}
\begin{document}
\maketitle
Write
\[
 F_n(T)=\sum_{j=0}^n\binom nj(6j)!T^j.
\]
All statements concern the ordinary factorial shape one. The
\href{MIXED_PRIME_RESIDUAL_WINDOWS.tex}{real-place proof} excludes
odd-degree rational square divisors, and the
\href{SLOPE_COVERAGE_AND_SHIFTED_WINDOWS.tex}{shifted-polygon proof}
retains the literal increasing-degree Newton convention. The following
user-provided arguments sharpen its large-degree range and extend its
finite residual kernels. Independent replay, formal verification,
external review and an all-index SFC conclusion are not claimed.

\section{An exposed primitive edge and the preceding-prime bound}
\begin{theorem}\label{thm:edge}
Let $q>6$ be prime, $n/2<q\le n$, and $q^2>6n$. Put $d=n-q$.
The Newton polygon of $F_n$ has an exposed edge of width $q$ and
height six. Its reduced slope denominator is $q$, and consequently
\[
 H^2\mid F_n,\quad H\in\Q_q[T]
 \quad\Longrightarrow\quad\deg H\le\lfloor d/2\rfloor. \tag{1}
\]
\end{theorem}
\begin{proof}
For $A_j=\binom nj(6j)!$, Legendre's formula and $n<2q$ give
\[
 v_q(A_j)=\left\lfloor\frac{6j}{q}\right\rfloor
             +\mathbf1_{\{d<j<q\}}. \tag{2}
\]
For $0\le h\le d$, the adjusted height $v_q(A_j)-6j/q$ at both
$j=h$ and $j=h+q$ is $-(6h\bmod q)/q$. At the remaining indices
$d<j<q$ it is $1-(6j\bmod q)/q>0$.
Since $d<q$ and $q>6$, multiplication by six is injective modulo $q$
on $0,\ldots,d$. There is a unique maximizer $h_*$ of $6h\bmod q$.
The supporting line of slope $6/q$ therefore touches exactly
$h_*,h_*+q$. This proves the exposed edge and its primitive denominator.
Newton polygons of products add edge widths at each slope. A square
divisor cannot occupy any part of that width-$q$ edge: its own positive
width there would be a multiple of $q$ and doubling would exceed $q$.
The remaining width of the whole polygon is $n-q=d$, so its doubled
degree fits into at most $d$. This proves (1).
\end{proof}

\begin{corollary}\label{cor:gap}
Let $q_n$ be the largest prime at most $n$, and $d_n=n-q_n$.
For every $n\ge30$ and every rational square divisor $H^2\mid F_n$,
\[
 \deg H\le\left\lfloor\frac{d_n}{2}\right\rfloor,
 \qquad \deg H\le2\left\lfloor\frac{d_n}{4}\right\rfloor. \tag{3}
\]
The second bound uses real-place parity.
\end{corollary}
\begin{proof}
The pinned \href{sources/nagura-1952-prime-interval.pdf}{Nagura theorem}
at $x=5n/6\ge25$ supplies a prime in $(5n/6,n)$, so $q_n>5n/6$.
For $n\ge30$ this implies $q_n>6$ and $q_n^2>6n$.
Apply Theorem~\ref{thm:edge}, then round down to an even degree.
\end{proof}
For $n=10^6$, $q_n=999983$ and $d_n=17$, so the only positive
remaining degrees under (3) are $2,4,6,8$.

The classical prime number theorem, in the precise form
\href{https://dlmf.nist.gov/27.2.E3}{DLMF 27.2.3}, implies $d_n=o(n)$:
for each fixed $\varepsilon>0$, the difference
$\pi(n)-\pi((1-\varepsilon)n)\sim\varepsilon n/\log n$ is eventually
positive. This asymptotic input is imported and credited. No numerical
cutoff for the later fixed-multiplier corollary is claimed.

\section{The explicit one-fourteenth bound}
\begin{theorem}\label{thm:fourteenth}
For every $n\ge30$, any rational square divisor of $F_n$ has
\[
 \boxed{\deg H<n/14.} \tag{4}
\]
\end{theorem}
\begin{proof}
Choose a prime $5n/6<q<n$ by Nagura and put $d=n-q$.
If $q>6n/7$, (1) gives $\deg H\le d/2<n/14$.
Otherwise put $L=\lfloor q/6\rfloor$. We have
$q/6\le d<q/5$ and $q\equiv1$ or $5\pmod6$.
Among $0\le h\le d$, the residue $6h\bmod q$ has its unique
maximum at $h=L$: there it is $q-(q\bmod6)$, while after the wrap
it is at most $6d-q<q/5$. Formula (2) now gives the exact vertices
\[
 (0,0),\quad(L,0),\quad(L+q,6),\quad(n,7). \tag{5}
\]
The middle edge has denominator $q$ and width $q$. The last edge
has height one, so its denominator equals its positive width $d-L$.
A doubled square divisor cannot use either edge. It therefore has
$\deg H\le L/2<q/12\le n/14$, proving (4).
\end{proof}

\begin{corollary}\label{cor:blocks}
The following are squarefree for all $r\ge1$:
\[
 \begin{array}{ll}
 F_{ap^r+i},&1\le a\le14,\quad0\le i\le3,\quad p>6\max(a,i),\\
 F_{ap^r+i},&1\le a\le27,\quad i=0,1,\quad p>6a,\\
 F_{28p^r},&p>168.
 \end{array} \tag{6}
\]
\end{corollary}
\begin{proof}
Put $P=p^r$. The retained shifted polygon forces
$\deg H=h+Pt$ with $h\le\lfloor i/2\rfloor$ and
$t\le\lfloor a/2\rfloor$. For $i\le3$, parity excludes a nonzero
purely horizontal contribution, and every other positive even degree
is at least $P+1$. If $a\le14$, then
$n/14=(aP+i)/14<P+1$; (4) excludes all such degrees.
For $i=0,1$, $h=0$, so the first positive even contribution is $2P$.
For $a\le27$, $n/14<2P$; at $a=28,i=0$, $n/14=2P$ and the
strict bound (4) still excludes it. If $n<30$, the prime inequality
forces $a\le2$ and these cases are already covered by the retained
shifted-parity and full-square proofs. No new bounded-index scan is
used for these universal conclusions.
\end{proof}

\section{A first-subresultant certificate for small residual kernels}
For $d\ge2$, define the integer data
\[
 R_d=\operatorname{Res}(F_d,F_d'),\qquad
 S_{d,1}=\operatorname{Sres}_1(F_d,F_d')=A_dT+B_d,
 \qquad C_d=\gcd(R_d,A_d,B_d)\ge0. \tag{7}
\]
The subresultant is used in its full integer normalization, with no
primitive-part or monic rescaling.

\begin{theorem}\label{thm:subres}
If $C_d>0$ and every prime divisor of $C_d$ is at most $6d$, then
for every prime $q>6d$,
\[
 \deg\gcd(\overline{F_d},\overline{F_d}')\le1. \tag{8}
\]
Every square divisor of that residual has degree at most one.
\end{theorem}
\begin{proof}
The leading coefficients of $F_d$ and $F_d'$ are units modulo $q$.
Resultants and subresultants are determinants over the integer
coefficient ring, so they specialize modulo $q$ without division.
Both $R_d$ and $S_{d,1}$ are integer polynomial combinations of
$F_d,F_d'$. If their reduced gcd had degree at least two, it would
divide the constant $R_d\bmod q$ and the polynomial $S_{d,1}\bmod q$
of degree at most one. Both must therefore vanish identically.
Then $q$ divides $R_d,A_d,B_d$, hence $C_d$, contradicting smoothness.
Finally, $K^2\mid\overline{F_d}$ implies
$K\mid\gcd(\overline{F_d},\overline{F_d}')$.
\end{proof}

For explicit division-free evidence for the linear combination,
form the columns $T^{d-3}F_d,\ldots,F_d$ and
$T^{d-2}F_d',\ldots,F_d'$. Take their coefficient matrix in degrees
$2d-3,\ldots,2$, and its signed maximal minors. These give an integer
null vector; its corresponding polynomial combination has degree
at most one and equals $S_{d,1}$ up to sign. For $d=2$ the matrix is
empty and the combination is simply $F_d'$.
The companion checker independently reconstructs these minors.

\href{small_kernel_certificates.json}{The integer certificate packet}
contains $R_d,A_d,B_d,C_d$, a Bezout identity for their integer gcd,
and a complete factorization of $C_d$ into primes at most $6d$,
for every $2\le d\le12$. All eleven remainders after removing these
small primes are one. In particular,
\[
 C_2=2^5 3^2 5,\qquad
 C_3=2^{30}3^{17}5^6 7^3 11^2 13\cdot17. \tag{9}
\]
This proves a finite list of \emph{universal prime} statements; it is
not a theorem for arbitrary $d$. Offsets zero and one are direct.

\begin{corollary}[Horizontal residual transfer]\label{cor:horizontal}
Let $n=q+d$ with $q>6\max(1,d)$ prime and $0\le d\le12$.
Then $F_n$ is squarefree. More generally,
\[
 \boxed{F_{p^r+i}\text{ is squarefree},\quad
 0\le i\le12,\quad p>6\max(1,i),\quad r\ge1.} \tag{10}
\]
\end{corollary}
\begin{proof}
In the first case, the Newton polygon has horizontal edge of length
$d$ followed by width $q$ and denominator $q$. Higher coefficients
vanish modulo $q$ and Lucas gives
$\overline{F_n}=\overline{F_d}$. A square divisor lies entirely on the
horizontal edge. After normalizing its constant coefficient to one,
its coefficients are $q$-integral and its leading coefficient is a
unit. Its unit leading coefficient also makes the quotient integral,
so reduction preserves its degree and square-divides $\overline{F_d}$. For $d\le12$, (8) bounds that degree by one, and
real-place parity excludes one. The cases $d=0,1$ are immediate.

For $n=p^r+i$, the retained shifted-polygon theorem has exactly the
same horizontal residual $\overline{F_i}$; its other edge has width
and denominator $p^r$, so a square divisor cannot occupy it. The
same residual certificate and parity finish for every $r$.
\end{proof}

\begin{corollary}[Combining both reusable results]\label{cor:rectangle}
For all $r\ge1$,
\[
 \boxed{F_{ap^r+i}\text{ is squarefree},\quad
 1\le a\le14,\quad0\le i\le12,\quad p>6\max(a,i).} \tag{11}
\]
For every fixed $a\ge1$ and fixed $0\le i\le12$, the same family
is squarefree for all sufficiently large $P=p^r$ satisfying the prime
condition. No effective cutoff is supplied for this eventual statement.
\end{corollary}
\begin{proof}
At horizontal slope, the weighted initial-product identity makes the
degree-$h$ residual of $H$ square-divide $\overline{F_i}$. Thus the
small-kernel certificate restricts $h$ to at most one, even when $i>3$. Parity then makes every positive
remaining contribution $h+Pt$ at least $P+1$.
For $a\le14$, $n/14\le P+12/14<P+1$, so (4) excludes it when
$n\ge30$. If $i\ge4$, the prime condition already gives $n\ge33$;
smaller indices therefore have $a\le2,i\le3$ and are covered as above.

For fixed $a,i$ without a uniform bound on $a$, use (3) and the
prime number theorem to obtain $\deg H=o(n)=o(P)$. Eventually no
positive multiple-of-$P$ contribution can occur, so $t=0$. The
remaining horizontal degree is at most one, and parity excludes one.
This is a qualitative existence argument for a cutoff, not its computation.
\end{proof}

\section{Executable boundary}
\href{prime_gap_windows.py}{The implementation} exports the exposed-edge,
preceding-prime and explicit one-fourteenth bounds, the retained eleven
integer kernel certificates and their near-prime/shifted consequences.
\href{verify_prime_gap_windows.py}{Exact controls} compare literal Newton
polygons and normalized horizontal residuals, reconstruct first
subresultants by determinant minors, and replay all integer gcd identities
and smoothness factorizations. The million-index example uses the prime-gap
formula, without constructing a million-degree moment polynomial.

All-index exponent-six SFC remains open. The remaining task concerns
small even degrees within the prime-gap bound and uniform residual-kernel
behavior beyond the retained prefix. The effective family (11), the
qualitative eventual statement and finite coefficient controls are
separate assertions. There is no Gamma-shape generalization of the new
ordinary factorial arguments or novelty claim.
\end{document}
